\documentclass[11pt]{article}

\usepackage{amsmath,amssymb}
\usepackage{booktabs}
\usepackage{xcolor}
\usepackage{xspace}
\usepackage[margin=1.4in]{geometry}
\usepackage[colorlinks=true,
            linkcolor=blue,
            urlcolor=blue]{hyperref}

\newcommand{\spx}{SPHINCS$^+$\xspace}

\title{A Fast Formula for Generic Security of SPHINCS$^+$}

\author{conduition \\ \href{https://brink.dev}{\emph{Brink Technology}} \\ \small \texttt{conduition@proton.me}}

\date{2026-11-11}

\begin{document}

\maketitle

\begin{abstract}
Generic attacks on SPHINCS$^+$ are feasible if a signer reveals a sufficient number of few-time signatures. Finding the exact probability of a successful generic forgery attempt typically requires computing a sum with a number of terms linear in the number of SPHINCS$^+$ signing oracle queries, which may be as high as $2^{64}$, and in practice even approximations require thousands of terms (or more) for sufficient accuracy on some parameter sets. In this note we derive an efficient formula for the generic security of any SPHINCS$^+$ parameter set, requiring only $k+1$ terms computed with high precision. We then generalize to acquire a closed-form formula for the generic security of any SPHINCS$^+$ instantiation whose few-time signature scheme's forgery security can be written as an exponential polynomial - This includes well-known schemes such as PORS+FP.
\end{abstract}

{\let\thefootnote\relax \footnotetext{This work is fully human-authored, made possible by a Brink grant.}}

\newpage
\section{Introduction}

We begin with a focus on the NIST-standardized \spx instance called SLH-DSA, with FORS as its few-time signature component \cite{sphincsplus}.

Let $h$ be the height of the \spx hypertree. Let $q_s$ be the number of signing queries. Let $k$ be the number of FORS trees, and $a$ their height.

\emph{Interleaved target subset resilience} (ITSR) is the security notion governing \spx forgery security with-respect to FORS key overuse. In \cite{sphincsplus} section 5, the authors derive a closed-form formula for ITSR security assuming the message hashing function is a random keyed function. We briefly review their formula.

Bernstein et al write $\mathrm{DarkSide}_\gamma$ for the probability that after exactly $\gamma$ signatures issued by a FORS key, a forgery attempt hits an exposed leaf in a FORS tree. Expanded, we have

\begin{equation}
  \mathrm{DarkSide}_\gamma := 1 - \left( 1 - \frac{1}{2^a} \right)^\gamma
\end{equation}

The overall \spx generic ITSR forgery probability $\sigma$ is then defined as a sum over the forgery probabilities for each possible number of revealed signatures.

\begin{align}
\label{eq:sigma-naive}
  \sigma &:= \sum_{\gamma=1}^{q_s} \mathrm{DarkSide}^k_\gamma \binom{q_s}{\gamma} \left( 1 - \frac{1}{2^h} \right)^{q_s - \gamma} \left( \frac{1}{2^{h}} \right)^\gamma \\
  &= \sum_{\gamma=1}^{q_s} \left( 1 - \left( 1 - \frac{1}{2^a} \right)^\gamma \right)^k \binom{q_s}{\gamma} \left( 1 - \frac{1}{2^h} \right)^{q_s - \gamma} \left( \frac{1}{2^{h}} \right)^\gamma
\end{align}

$\lambda = -\log_2 \sigma$ is precisely the ``bit security" level of \spx against generic FORS forgery attacks.

\section{Motivation}

Computing $\sigma$ precisely requires evaluating the sum in (\ref{eq:sigma-naive}) over $q_s$ terms, which for many parameter sets is exceedingly large, e.g. $q_s = 2^{64}$ for NIST's SLH-DSA parameter sets \cite{fips205}.

Most of the weight in this probability sum sits on the first few terms, so $\sigma$ is typically \emph{approximated} by summing only the leading terms until the contributions of additional terms fall below some threshold. This is imprecise, and does not scale when $q_s \gg 2^h$.

\cite{smaller_sphincs} (Section 5) gives an approximation of (\ref{eq:sigma-naive}) in terms of a Poisson distribution which can be more efficiently computed per term, but it is still an approximation which deviates from the true $\sigma$ when $q_s \ll 2^h$.

In the following section, we derive an \emph{exact} closed-form formula for $\sigma$, requiring only $O(k)$ computations.

\section{FORS}
\label{sec:fors}

Expand $\mathrm{DarkSide}^k_\gamma$ with the binomial theorem identity $(x + y)^n = \sum_{i=0}^n \binom{n}{i} x^i y^{n-i}$.

\begin{align}
  \mathrm{DarkSide}^k_\gamma
    &= \left( 1 - \left( 1 - \frac{1}{2^a} \right)^\gamma \right)^k \\
    &= \sum_{i=0}^k \binom{k}{i} (-1)^i \left( 1 - \frac{1}{2^a} \right)^{i \gamma}
\end{align}

Let $P_\text{select}(\gamma) := \binom{q_s}{\gamma} \left( 1 - \frac{1}{2^h} \right)^{q_s - \gamma} \left( \frac{1}{2^h} \right)^\gamma$ be the binomial probability of exactly $\gamma$ successes after $q_s$ trials, with a per-trial success rate of $\frac{1}{2^h}$.

Substitute into (\ref{eq:sigma-naive}).

\begin{align}
  \sigma &= \sum_{\gamma=1}^{q_s} \mathrm{DarkSide}^k_\gamma \binom{q_s}{\gamma} \left( 1 - \frac{1}{2^h} \right)^{q_s - \gamma} \left( \frac{1}{2^{h}} \right)^\gamma \\
         &= \sum_{\gamma=1}^{q_s} \sum_{i=0}^k \binom{k}{i} (-1)^i \left( 1 - \frac{1}{2^a} \right)^{i \gamma} P_\text{select}(\gamma)
\end{align}

Since $\mathrm{DarkSide}_0 = 0$, we may extend the sum to include a $\gamma = 0$ term which contributes nothing\footnote{Intuitively: With zero preimages exposed, an attacker cannot hit any of them.}. Then swap the sums and factor out.

\begin{align}
  \sigma &= \sum_{\gamma=0}^{q_s} \sum_{i=0}^k \binom{k}{i} (-1)^i \left( 1 - \frac{1}{2^a} \right)^{i \gamma} P_\text{select}(\gamma)\\
         &= \sum_{i=0}^k \sum_{\gamma=0}^{q_s} \binom{k}{i} (-1)^i \left( 1 - \frac{1}{2^a} \right)^{i \gamma} P_\text{select}(\gamma) \\
         &= \sum_{i=0}^k \binom{k}{i} (-1)^i \sum_{\gamma=0}^{q_s} \left[ \left( 1 - \frac{1}{2^a} \right)^i \right]^\gamma P_\text{select}(\gamma)
\end{align}

Let $t(i) = \left( 1 - \frac{1}{2^a} \right)^i$.

\begin{equation}
  \sigma = \sum_{i=0}^k \binom{k}{i} (-1)^i \sum_{\gamma=0}^{q_s} t(i)^\gamma \cdot P_\text{select}(\gamma)
\end{equation}

Note that $P_\text{select}$ is a binomial probability mass function, so the inner sum is a \emph{probability generating function} (PGF) giving the expected value of $t(i)^X$ for random variable $X \sim \mathrm{Binomial} \left( q_s, \frac{1}{2^h} \right)$.

\begin{equation}
  \sigma = \sum_{i=0}^k \binom{k}{i} (-1)^i \cdot \mathbb E \left[ t(i)^X \right]
\end{equation}

The binomial PGF has a closed form, again using the binomial theorem identity.

\begin{align}
  \mathbb E \left[ t(i)^X \right]
    &= \sum_{\gamma=0}^{q_s} t(i)^\gamma \cdot P_\text{select}(\gamma) \\
    &= \sum_{\gamma=0}^{q_s} t(i)^\gamma \binom{q_s}{\gamma} \left( 1 - \frac{1}{2^h} \right)^{q_s - \gamma} \left( \frac{1}{2^h} \right)^\gamma \\
    &= \sum_{\gamma=0}^{q_s} \binom{q_s}{\gamma} \left( \frac{t(i)}{2^h} \right)^\gamma \left( 1 - \frac{1}{2^h} \right)^{q_s - \gamma}  \\
    &= \left( \frac{t(i)}{2^h} + \left(1 - \frac{1}{2^h} \right) \right)^{q_s} \\
    &= \left(1 +  \frac{t(i) - 1}{2^h} \right)^{q_s} \\
    &= \left(1 - \frac{1 - t(i)}{2^h} \right)^{q_s} \\
    &= \left(1 - \frac{1 - \left( 1 - \frac{1}{2^a} \right)^i}{2^h} \right)^{q_s}
\end{align}

This gives a closed-form solution for $\sigma$ with only $k+1$ terms.

\begin{equation}
\label{eq:fors-closed}.
\boxed{
  \sigma = \sum_{i=0}^k \binom{k}{i} (-1)^i \cdot \left(1 - \frac{1 - \left( 1 - \frac{1}{2^a} \right)^i}{2^h} \right)^{q_s}
}
\end{equation}

This offers an efficient formula for the exact \spx ITSR bit-security level against generic attacks: $\lambda = - \log_2 \sigma$.

We note the sum experiences catastrophic cancellation, and so requires numeric precision of at least $\log_2 \left( \frac{\binom{k}{\lfloor k/2 \rfloor}}{\sigma} \right)$ bits.

\subsection{Implementation}

We provide a supplemental implementation of the exact formula from (\ref{eq:fors-closed}) in Python, and empirically compare our .formula against the approximations from \cite{sphincsplus, smaller_sphincs} on various \spx parameter sets. It is available at

\url{https://github.com/
  &conduition/sphincs-itsr-formula-note}

\section{Generalization}

While FORS is the FTS scheme standardized for \spx in FIPS-205, novel FTS schemes have since been proposed.

In this section, we generalize the derivation from Section \ref{sec:fors} so that it may be applied to (claimed) security bounds in proposed FTS schemes such as PORS+FP \cite{forced_pruning} or BPORS+FP \cite{bucket_thinning}.

Let $P_F(\gamma)$ be a function that returns the forgery probability of an FTS scheme $F$ after $\gamma$ signatures released by a key. When scheme $F$ is used in \spx, the upper bound for the generic forgery probability $\sigma_F$ of a public key after releasing $q_s$ signatures is

\begin{align}
  \sigma_F
    &\le \sum_{\gamma=0}^{q_s} \binom{q_s}{\gamma} \left( 1 - \frac{1}{2^h} \right)^{q_s - \gamma} \left( \frac{1}{2^h} \right)^\gamma P_F(\gamma) \\
    &= \sum_{\gamma=0}^{q_s} P_\text{select}(\gamma) \cdot P_F(\gamma)
\end{align}

Recalling that $P_\text{select}$ is the binomial PMF:

\begin{equation}
  X \sim \mathrm{Binomial} \left( q_s, \frac{1}{2^h} \right)
\end{equation}
\begin{equation}
  \sigma_F \le \mathbb E \left[ P_F(X) \right]
\end{equation}

For convenience, let $p = \frac{1}{2^h}$ and let $\mathcal W(g) := \mathbb E \left[ g(X) \right]$ for $X \sim \mathrm{Binomial} \left( q_s, p \right)$.

\begin{align}
  \mathcal W(g)
    &:= \sum_{\gamma=0}^{q_s} P_\text{select}(\gamma) \cdot g(\gamma) \\
    &= \sum_{\gamma=0}^{q_s} \binom{q_s}{\gamma} \left( 1 - p \right)^{q_s - \gamma} p^\gamma g(\gamma)
    \label{eq:exp-op}
\end{align}

\paragraph{$\mathcal W$ is linear.} This follows from the fact that the expectation of sums is the sum of expectations.

\begin{equation}
  \mathbb E[a(X) + b(X)] = \mathbb E[a(X)] + \mathbb E[b(X)]
\end{equation}
\begin{equation}
  \mathcal W(a + b) = \mathcal W(a) + \mathcal W(b)
\end{equation}

To compute $\mathcal W(g)$, we break $g$ down into $M$ additive pieces $\{g_1, g_2, \ldots g_M\}$ such that

\begin{equation}
  g(\gamma) = \sum_{i=1}^M g_i(\gamma)
\end{equation}

Then we will have

\begin{equation}
  \mathcal W(g) = \sum_{i=1}^M \mathcal W(g_i)
\end{equation}

% \paragraph{$\mathcal W(g)$ has identities for certain functions $g$.} As shown earlier (TODO secref), the binomial expectation $\mathbb E \left[ z^X \right]$ for any $z \in \mathbb R$ is one such example.

\bigskip

We will show that there exists a broad class of functions for which this form of $\mathcal W(g)$ is efficiently computable, giving an upper bound for $\sigma_F$. Namely, we restrict $g$ to be an exponential polynomial

\begin{equation}
\label{eq:exponential-polynomial}
  g(\gamma) := \sum_{i=1}^M f_i(\gamma) \cdot z_i^\gamma
\end{equation}

where $\{f_1, f_2, \ldots f_M\}$ is a set of real polynomials, and $\{z_1, z_2, \ldots z_M\}$ is a set of exponent bases.

Let $d_i := \deg f_i$.

Write each polynomial $f_i(\gamma)$ using a binomial-coefficient basis,\footnote{Any $d$-degree polynomial $f(x) = \sum_{i=0}^d a_i x^i$ can be rewritten in binomial-coefficient basis, as $$f(x) = \sum_{i=0}^d \binom{x}{i} i! \sum_{j=i}^d a_j \cdot S(j, i)$$ where $S(j, i)$ denotes Stirling numbers of the second kind.} with coefficients $\{c_{i, 0} \ldots c_{i, d_i}\}$.

\begin{equation}
  f_i(\gamma) = \sum_{j=0}^{d_i} c_{i,j} \binom{\gamma}{j}
\end{equation}

We then have $M$ components which sum to $g$.

\begin{equation}
  g_i(\gamma) := f_i(\gamma) \cdot z_i^\gamma
\end{equation}
\begin{equation}
  g(\gamma) = \sum_{i=1}^M g_i(\gamma)
\end{equation}

Recall that by linearity of $\mathcal W$, we have

\begin{equation}
\label{eq:W-of-g}
  \mathcal W(g) = \sum_{i=1}^M \mathcal W(g_i)
\end{equation}

Focus on a single term $\mathcal W(g_i) = \mathbb E \left[ \sum_{j=0}^{d_i} c_{i,j} \binom{X}{j} \cdot z_i^X \right]$. Note none of the sum indices depend on the random variable $X$, and nor do the coefficients $\{c_{i,j}\}$. By linearity,

\begin{equation}
   \mathbb E \left[ \sum_{j=0}^{d_i} c_{i,j} \binom{X}{j} \cdot z_i^X \right]
    = \sum_{j=0}^{d_i} c_{i,j} \cdot \mathbb E \left[ \binom{X}{j} \cdot z_i^X \right]
\end{equation}

We further refine our focus on computing $\mathcal W \left( \left[ \binom{X}{j} \cdot z^X \right] \right)$ for $j \in \mathbb N$, and $z \in \mathbb R$. Our goal will be to rearrange this definition into a form we can factor using the binomial theorem.

Expand (\ref{eq:exp-op}).

\begin{align}
  \mathbb E \left[ \binom{X}{j} \cdot z_i^X \right]
    &= \sum_{\gamma=0}^{q_s} \binom{q_s}{\gamma} (1 - p)^{q_s - \gamma} p^\gamma \left( \binom{\gamma}{j} \cdot z^\gamma \right) \\
    &= \sum_{\gamma=0}^{q_s} \binom{q_s}{\gamma} \binom{\gamma}{j} (1 - p)^{q_s - \gamma} (pz)^\gamma
\end{align}

Use the binomial identity $\binom{k}{n} \binom{n}{m} = \binom{k}{m} \binom{k-m}{n-m}$, and factor out the $\gamma$-free part.

\begin{align}
    &= \sum_{\gamma=0}^{q_s} \binom{q_s}{j} \binom{q_s - j}{\gamma - j} (1 - p)^{q_s - \gamma} (pz)^\gamma \\
    &= \binom{q_s}{j} \sum_{\gamma=0}^{q_s} \binom{q_s - j}{\gamma - j} (1 - p)^{q_s - \gamma} (pz)^\gamma
\end{align}

Extract a factor $(pz)^j$.

\begin{align}
  &= (pz)^j \binom{q_s}{j} \sum_{\gamma=0}^{q_s} \binom{q_s - j}{\gamma - j} (1 - p)^{q_s - \gamma} (pz)^{\gamma - j}
\end{align}

Terms with $\gamma < j$ contribute nothing (the $\binom{q_s - j}{\gamma - j}$ factor is zero), so we may drop the first $j$ terms, and then shift the index by $j$ to clean up.

\begin{align}
    &= (pz)^j \binom{q_s}{j} \sum_{\gamma=j}^{q_s} \binom{q_s - j}{\gamma - j} (1 - p)^{q_s - \gamma} (pz)^{\gamma - j} \\
    &= (pz)^j \binom{q_s}{j} \sum_{\gamma=0}^{q_s - j} \binom{q_s - j}{\gamma} (1 - p)^{q_s - j - \gamma} (pz)^{\gamma}
\end{align}

We may now apply the binomial theorem to factor the sum,

\begin{equation}
  \sum_{\gamma=0}^N \binom{N}{\gamma} x^\gamma y^{N-\gamma} = (x + y)^N
\end{equation}

by substituting $N = q_s - j$; $x = pz$; $y = 1 - p$.

\begin{align}
  \mathbb E \left[ \binom{X}{j} \cdot z_i^X \right]
    &= (pz)^j \binom{q_s}{j} \left(1 - p + pz \right)^{q_s - j}
\end{align}

Finally we can reconnect this thread of work to fully evaluate $\mathcal W$ on the exponential polynomial $g$ as in (\ref{eq:W-of-g}).

\begin{align}
  \mathcal W(g)
    &= \sum_{i = 1}^M \mathcal W(g_i) \\
    &= \sum_{i = 1}^M \sum_{j=0}^{d_i} c_{i, j} \cdot \mathbb E \left[ \binom{X}{j} \cdot z_i^X \right]
\end{align}

\begin{equation}
\label{eq:generalization}
  \boxed{
    \mathcal W(g) = \sum_{i = 1}^M \sum_{j=0}^{d_i} c_{i, j} (p z_i)^j \binom{q_s}{j} \left(1 - p + p z_i \right)^{q_s - j}
  }
\end{equation}

This is an exact and closed-form formula to compute $\mathbb E [g(X)]$ for $X \sim \mathrm{Binomial}(q_s, p)$ with $\sum_{i=1}^M (d_i + 1)$ terms. Contrast this with (\ref{eq:exp-op}), which requires $q_s$ terms to compute.

\subsection{Application}

If an FTS forgery probability curve $P_F(\gamma)$ can be expressed as an exponential polynomial, as in (\ref{eq:exponential-polynomial}), then $\mathcal W(P_F)$ provides a closed-form upper bound for $\sigma_F$.

\subsubsection{FORS}

Applied to FORS, we have

\begin{equation}
  P_F(\gamma) = \mathrm{DarkSide}_\gamma^k = \sum_{i=0}^k \binom{k}{i} (-1)^i \left( 1 - \frac{1}{2^a} \right)^{i \gamma}
\end{equation}

Written as an exponential polynomial, as in (\ref{eq:exponential-polynomial}), we have

\begin{equation}
  f_i(\gamma) = \binom{k}{i} (-1)^i
    \quad \quad
  z_i = \left( 1 - \frac{1}{2^a} \right)^i
\end{equation}
\begin{equation}
  P_F(\gamma) = \sum_{i=0}^k f_i(\gamma) \cdot z_i^\gamma
\end{equation}

Since every $f_i(\gamma)$ is just a constant, rewriting $f_i$ with binomial-coefficient basis gives

\begin{equation}
  f_i(\gamma) = \overbrace{\binom{\gamma}{0}}^1 \overbrace{\binom{k}{i} (-1)^i}^{c_{i, 0}}
\end{equation}

i.e. $c_{i, 0} = \binom{k}{i} (-1)^i$, while all $c_{i, j} = 0$ when $j > 0$.

Substituting into (\ref{eq:generalization}) with $M = k + 1$ (indexing from zero instead), we recover (\ref{eq:fors-closed}), the formula for generic security of \spx with FORS.

\begin{align}
  \mathcal W(P_F)
    &= \sum_{i = 0}^k \sum_{j=0}^{d_i} c_{i, j} (p z_i)^j \binom{q_s}{j} \left(1 - p + p z_i \right)^{q_s - j} \\
    &= \sum_{i = 0}^k c_{i, 0} \left(1 - p + p z_i \right)^{q_s} \\
    &= \sum_{i = 0}^k \binom{k}{i} (-1)^i \left(1 - (p - p z_i) \right)^{q_s} \\
    &= \sum_{i = 0}^k \binom{k}{i} (-1)^i \left(1 - p \left( 1 - z_i \right) \right)^{q_s} \\
    &= \sum_{i = 0}^k \binom{k}{i} (-1)^i \left(1 - \frac{1 - \left( 1 - \frac{1}{2^a} \right)^i}{2^h} \right)^{q_s} \\
    &= \sigma_F
\end{align}

\subsubsection{PORS+FP}

Applied to PORS+FP \cite{forced_pruning}, with fixed parameters $k$ and $t$, we have

\begin{equation}
\label{eq:pors-pf}
  P_F(\gamma) = \binom{k \gamma}{k} / \binom{t}{k}
\end{equation}

This is already written as a polynomial in $\gamma$ of degree $k$, with $M = 1$, and all bases $z_i = 1$. Thus we drop the summation over $i$.

\begin{align}
  \mathcal W(P_F)
    &= \sum_{i = 1}^M \sum_{j=0}^{d_i} c_{i, j} (p z_i)^j \binom{q_s}{j} \left(1 - p + p z_i \right)^{q_s - j} \\
    &= \sum_{j=0}^k c_j\ p^j \binom{q_s}{j} \left(1 - p + p \right)^{q_s - j} \\
    &= \sum_{j=0}^k c_j\ p^j \binom{q_s}{j} \label{eq:pors-forge-succinct}
\end{align}

It remains to find the binomial-coefficient basis coefficients $\{c_0, c_1, \ldots, c_k\}$ such that $P_F(\gamma) = \sum_{j=0}^k c_j \binom{\gamma}{j} = \binom{k \gamma}{k} / \binom{t}{k}$. We will construct an apparatus to compute these coefficients using only the given form of $P_F$.

Let $\Delta$ denote the \emph{forward difference} operator.

\begin{equation}
\label{eq:forward-diff}
  (\Delta f)(\gamma) = f(\gamma + 1) - f(\gamma)
\end{equation}

Let $\Delta^j f$ denote the $j$-th iteration of $\Delta$ on $f$.

\begin{align}
  (\Delta^j f)(\gamma) &= \left(\Delta (\Delta^{j-1} f) \right)(\gamma) \\
    &= (\Delta^{j-1} f)(\gamma + 1) - (\Delta^{j-1} f)(\gamma)
\end{align}

Apply $\Delta$ to the exponential polynomial form of $P_F$.

\begin{align}
  P_F(\gamma) &= \sum_{j=0}^k c_j \binom{\gamma}{j} \\
  \Delta P_F(\gamma) &= \sum_{j=0}^k c_j \binom{\gamma + 1}{j} - \sum_{j=0}^k c_j \binom{\gamma}{j} \\
  \Delta P_F(\gamma) &= \sum_{j=0}^k c_j \left( \binom{\gamma + 1}{j} - \binom{\gamma}{j} \right)
\end{align}

Invoke Pascal's rule $\binom{n}{m} = \binom{n - 1}{m} + \binom{n - 1}{m - 1}$ and apply $\Delta$ repeatedly.

\begin{align}
  \Delta P_F(\gamma) &= \sum_{j=1}^k c_j \binom{\gamma}{j - 1} \\
  \Delta^\ell P_F(\gamma) &= \sum_{j=\ell}^k c_j \binom{\gamma}{j - \ell}
\end{align}

Note the first $\ell$ terms are dropped because each contains a factor $\binom{\gamma}{j - \ell} = 0$ if $j < \ell$. If we evaluate $\Delta^\ell P_F(0)$, the trailing terms contain a factor $\binom{0}{j - \ell} = 0$ if $j > \ell$. This leaves only the $j = \ell$ term non-zero, for which we have

\begin{align}
  \Delta^\ell P_F(0)
    &= \sum_{j=\ell}^k c_j \binom{0}{j - \ell} \\
    &= c_\ell \binom{0}{0} \\
    &= c_\ell
\end{align}

Concretely: \textbf{Evaluating $\Delta^\ell P_F(0)$ gives the coefficient $c_\ell$.} In practice one could compute $\Delta^\ell P_F(0)$ using only the forward difference definition from (\ref{eq:forward-diff}) by building a table of evaluations of $P_F$, but we seek a closed form to compute the coefficients analytically.

Let $(S f)(\gamma) := f(\gamma + 1)$ be the shift operator, and let $(I f)(\gamma) = f(\gamma)$ be the identity operator. We can write $\Delta = S - I$, and $\Delta^j = (S - I)^j$.

$S$ and $I$ commute as operators under composition ($I$ is trivially commutative). By the binomial identity for commutative operators,

\begin{align}
  \Delta^\ell f
    &= (S - I)^\ell f \\
    &= \sum_{j = 0}^\ell \binom{\ell}{j} (-1)^{\ell - j} S^j f
\end{align}

Because $S^j f(0) = f(j)$, we can compute the coefficients using $P_F$.

\begin{align}
  c_\ell
    &= \Delta^\ell P_F(0) \\
    &= \sum_{j = 0}^\ell \binom{\ell}{j} (-1)^{\ell - j} (S^j P_F)(0) \\
    &= \sum_{j = 0}^\ell \binom{\ell}{j} (-1)^{\ell - j} P_F(j)
\end{align}

Substituting (\ref{eq:pors-pf}),

\begin{align}
  c_\ell
    &= \sum_{j = 0}^\ell \binom{\ell}{j} (-1)^{\ell - j} \binom{k j}{k} / \binom{t}{k} \\
    &= \binom{t}{k}^{-1} \sum_{j = 0}^\ell (-1)^{\ell - j} \binom{\ell}{j} \binom{k j}{k}
\end{align}

Finally substituting back into (\ref{eq:pors-forge-succinct}), we obtain an exact formula to compute Abri and Katz's upper bound for the generic security of PORS+FP in only $(k+1)(k+2)/2 = O(k^2)$ terms, rather than the $q_s$ terms in the bound formula (Equation 2) of \cite{forced_pruning}, which was derived from that of \cite{sphincsplus}.

\begin{equation}
\boxed{
  \mathcal W(P_F) = \binom{t}{k}^{-1} \sum_{\ell=0}^k p^\ell \binom{q_s}{\ell} \sum_{j = 0}^\ell (-1)^{\ell - j} \binom{\ell}{j} \binom{k j}{k}
}
\end{equation}

\begin{equation}
  \sigma_F \le \mathcal W(P_F)
\end{equation}

\newpage
\bibliography{references}
\bibliographystyle{alpha}

\end{document}
